\section{Algebraic geometry}
\source{fga-explained:page-0017}\source{fga-explained:page-0018}\source{fga-explained:page-0019}\source{fga-explained:page-0020}\source{fga-explained:page-0021}\source{fga-explained:page-0022}

All rings and algebras are commutative.  A finitely presented algebra is a
quotient of a polynomial algebra by a finitely generated ideal.

\begin{definition}[locally of finite presentation]\label{fga-def-1-1}\dcref{1.1}\source{fga-explained:page-0017}\uses{}
A morphism of schemes $f:X\to Y$ is locally of finite presentation if every
$x\in X$ has affine neighborhoods $U$ and $V$ of $x$ and $f(x)$, with
$f(U)\subseteq V$, such that $\OO_X(U)$ is a finitely presented
$\OO_Y(V)$-algebra.
\end{definition}\begin{proposition}\label{fga-prop-1-2}\dcref{1.2}\source{fga-explained:page-0017}\uses{fga-def-1-1}
Locally finitely presented morphisms are stable under restriction to suitable
affines, composition, and arbitrary base change.
\end{proposition}
\begin{proof}
The affine assertion is the definition.  Finite presentations compose by
substituting finitely many generators and relations, and tensoring a finite
presentation with the coordinate algebra of the base change gives a finite
presentation after pullback.
\end{proof}

\begin{definition}[quasi-compact morphism]\label{fga-def-1-3}\dcref{1.3}\source{fga-explained:page-0017}\uses{}
A morphism $f:X\to Y$ is quasi-compact if the inverse image of every
quasi-compact open of $Y$ is quasi-compact.
\end{definition}\begin{proposition}\label{fga-prop-1-4}\dcref{1.4}\source{fga-explained:page-0018}\uses{fga-def-1-3}
For a morphism $f:X\to Y$ the following are equivalent: $f$ is
quasi-compact; the inverse image of every affine open of $Y$ is
quasi-compact; and there is an affine open cover $Y=\bigcup V_i$ for which
$f^{-1}(V_i)$ is quasi-compact.  A morphism from a quasi-compact scheme to an
affine scheme is quasi-compact.
\end{proposition}
\begin{proof}
Use that a quasi-compact scheme is a finite union of affine opens.  Refining
an affine cover and taking finite subcovers proves the three implications.
\end{proof}

\begin{proposition}\label{fga-prop-1-6}\dcref{1.6}\source{fga-explained:page-0018}\uses{fga-def-1-3,fga-prop-1-4}
Quasi-compact morphisms are stable under composition and base change.
\end{proposition}
\begin{proof}
Apply the affine-cover criterion of \cref{fga-prop-1-4} successively to a
cover of the target and to its inverse image.
\end{proof}

\begin{definition}[flat and faithfully flat]\label{fga-def-1-7}\dcref{1.7}\source{fga-explained:page-0018}\source{fga-explained:page-0019}\uses{}
A morphism $f:X\to Y$ is flat when each local ring $\OO_{X,x}$ is flat over
$\OO_{Y,f(x)}$.  It is faithfully flat when it is flat and surjective.
\end{definition}\begin{proposition}\label{fga-prop-1-8}\dcref{1.8}\source{fga-explained:page-0018}\uses{fga-def-1-7}
Flatness is equivalent to affine flatness on every pair of compatible affine
neighborhoods.  In particular it can be checked on affine charts.
\end{proposition}
\begin{proof}
Localization of a flat module is flat.  Conversely, affine-local flatness
localizes to the stalks, giving the defining condition.
\end{proof}

\begin{proposition}\label{fga-prop-1-9}\dcref{1.9}\source{fga-explained:page-0018}\uses{fga-def-1-7}
Flat morphisms are stable under composition and arbitrary base change.
\end{proposition}
\begin{proof}
The tensor product of flat modules is flat, and base change is tensor product
on affine charts.
\end{proof}

\begin{proposition}\label{fga-prop-1-11}\dcref{1.11}\source{fga-explained:page-0019}\uses{fga-def-1-7}
For an $A$-algebra $B$, faithful flatness is equivalent to preservation and
reflection of exactness after tensoring with $B$; equivalently, a map of
$A$-modules is injective exactly when its base change is injective.  It is also
equivalent to flatness together with $MB=0\Rightarrow M=0$ for every
$A$-module $M$, or with $\mathfrak mB\ne B$ for every maximal ideal
$\mathfrak m\subset A$.
\end{proposition}
\begin{proof}
Tensoring preserves exactness under flatness.  Reflection follows from the
vanishing criterion, and the maximal-ideal formulation is the usual local
criterion for faithful flatness.
\end{proof}

\begin{proposition}\label{fga-prop-1-12}\dcref{1.12}\source{fga-explained:page-0019}\uses{fga-def-1-1,fga-def-1-7}
A flat morphism locally of finite presentation is open.
\end{proposition}

\begin{proposition}\label{fga-prop-1-13}\dcref{1.13}\source{fga-explained:page-0019}\uses{fga-prop-1-11,fga-prop-1-12}
If $f:X\to Y$ is faithfully flat, surjective and quasi-compact, a subset
$A\subseteq Y$ is open if and only if $f^{-1}(A)$ is open.
\end{proposition}
\begin{proof}
The forward implication is immediate.  For the converse, cover $Y$ by
affines which are images of quasi-compact opens of $X$ and use the affine
criterion for openness on each such piece.
\end{proof}

\begin{proposition}\label{fga-prop-1-15}\dcref{1.15}\source{fga-explained:page-0020}\uses{fga-prop-1-13}
In a cartesian square where $Y'\to Y$ is faithfully flat and either
quasi-compact or locally of finite presentation, each of the properties
separated, quasi-compact, locally of finite presentation, proper, affine,
finite, flat, smooth, unramified, etale, embedding, and closed embedding
descends from $X'\to Y'$ to $X\to Y$.
\end{proposition}
\begin{proof}
All listed properties are local on the target in the Zariski topology and are
preserved by base change.  Faithfully flat descent reduces the assertion to
the affine-local criteria.
\end{proof}

\section{Category theory}
\source{fga-explained:page-0020}\source{fga-explained:page-0021}\source{fga-explained:page-0022}

We use comma categories $(\mathcal C/X)$, opposite categories, fiber products,
and the projections $\operatorname{pr}_i$ with their evident functoriality.
A functor is fully faithful when it induces bijections on all hom-sets and
essentially surjective when every target object is isomorphic to an image.

\begin{proposition}\label{fga-prop-1-16}\dcref{1.16}\source{fga-explained:page-0020}\uses{fga-def-1-1}
A functor is an equivalence of categories if and only if it is fully faithful
and essentially surjective.
\end{proposition}
\begin{proof}
An inverse up to natural isomorphism gives both properties.  Conversely, choose
for each target object a preimage and transport morphisms through the hom-set
bijections; the choices assemble into a quasi-inverse.
\end{proof}

\begin{definition}[groupoid and product-preserving functor]\label{fga-def-1-17}\source{fga-explained:page-0021}\source{fga-explained:page-0022}\uses{}
A groupoid is a category in which every arrow is invertible.  A functor between
categories with finite products preserves finite products when it preserves a
terminal object and binary products (equivalently, all finite products).
\end{definition}